\section{Introduction}
Analyzing deployed JavaScript often requires understanding which library functions the delivered code implements. A security analyst, for example, may need to connect a bundled call target to a known library operation before interpreting its role in the application. Yet applications distribute library and application code together, with names and structure altered by bundling and minification. Recognizing that a package is present provides useful context, but leaves the correspondence between individual bundled functions and library implementations unresolved.

Existing techniques address parts of this problem. DEBUN uses property-order-graph (POG) fingerprints to identify libraries and versions~\cite{kim2025debun}; functions with little distinguishing structure remain difficult to identify independently. D-Bundlr recovers module and API information for downstream security analysis~\cite{xu2026dbundlr}. These advances motivate a broader recovery objective: identify library definitions, including short helpers, internal functions, and anonymous callbacks, across different bundled representations. Section~\ref{sec:background} explains the relevant transformations and the precise limitations of these techniques.

Our central observation is that a function need not be distinctive in isolation to be recoverable. An established correspondence can narrow the search for related functions: neighbors of the matched reference function provide candidates for neighbors of its bundled counterpart. Realizing this idea requires exposing relationships on both sides despite bundling, and assessing competing candidates without treating structural proximity as an identity guarantee.

We present ReCall, a method for recovering library function identities from supplied JavaScript files and candidate reference libraries. Code captured during webpage visits is one intended input setting. File acquisition and reference selection precede recovery; the method operates within the available code. ReCall's contributions are:
\begin{itemize}
\item \textbf{Iterative recovery beyond fingerprint anchors.} ReCall combines initial POG matches with neighborhood search, code-embedding ranking, and heuristic model validation. Accepted correspondences support further expansion to functions that lack distinctive fingerprints, without requiring a handwritten behavioral model for each target.
\item \textbf{Bundler-aware context for candidate generation.} A shared module and function-value representation supports independently constructed reference and bundled relationship graphs. Adapters normalize multiple loading and export protocols, allowing the search to cross boundaries that ordinary lexical relationships leave unresolved.
\item \textbf{An evaluation separating recovery decisions from graph reachability.} Model-driven experiments assess accepted identities and error propagation, while controlled connectivity simulations isolate whether graph construction exposes the correct candidates. This separation identifies both the potential of broader context and the remaining uncertainty in final matching accuracy.
\end{itemize}

The completed evidence is preliminary. With known references, the model-driven jQuery~4.0.0 pilot correctly recovers 1538/1629 source-map-supported bundled occurrences (94.41\%) across three minified builds, with 14 wrong identities and 45 unassessed promotions. This experiment uses the original local graph. Separately, on five previously untested repositories, the unchanged extended analyzer increases within-chunk reachability from 22,494/28,780 (78.16\%) to 26,402/28,780 supported pairs (91.74\%). The latter experiment uses oracle-confirmed seeds and ideal candidate selection; it measures graph reachability, not model recovery accuracy. Source-map labels serve evaluation only. Broader evaluation across bundlers and libraries/frameworks, validated transformation analysis, model-driven ablations, and configuration studies remain pending. The practical-use protocol remains to be defined.

The evaluation is organized around performance, ablation, tradeoff, and practical use; the last category remains an explicit placeholder.
